\chapter{Clasificación UNESCO del proyecto}
\label{anexo:codigos_unesco}

Este anexo recoge los códigos UNESCO asociados al proyecto TierraViva, tomando como referencia la Nomenclatura Internacional de UNESCO para los campos de Ciencia y Tecnología \cite{unesco_nomenclatura_cyt}. La clasificación se ha realizado a partir del alcance real del prototipo desarrollado: estaciones de campo autónomas para monitorización ambiental y de suelo, decisión local de riego, comunicación LTE/HTTP con ThingSpeak, almacenamiento histórico en Raspberry Pi, API y \textit{dashboard} web de supervisión.

La selección identifica los campos de conocimiento que describen de forma directa el trabajo realizado, diferenciando entre códigos principales, códigos secundarios y códigos complementarios u opcionales.

\section{Códigos principales}

Los códigos principales representan el núcleo técnico y funcional de TierraViva: riego automatizado, control local, adquisición de datos mediante sensores, comunicación remota y supervisión software.

\begingroup
\small
\renewcommand{\arraystretch}{1.18}
\rowcolors{2}{TableRowSoft}{white}

\begin{longtable}{|
>{\raggedright\arraybackslash}p{0.18\textwidth}|
>{\raggedright\arraybackslash}p{0.30\textwidth}|
>{\raggedright\arraybackslash}p{0.42\textwidth}|}

\caption{Códigos UNESCO principales asociados al proyecto TierraViva.}
\label{cuadro:codigos_unesco_principales}\\

\hline
\rowcolor{URJCRed}
\TVHead{Código UNESCO} & \TVHead{Denominación} & \TVHead{Justificación en TierraViva} \\
\hline
\endfirsthead

\hline
\rowcolor{URJCRed}
\TVHead{Código UNESCO} & \TVHead{Denominación} & \TVHead{Justificación en TierraViva} \\
\hline
\endhead

\hline
\multicolumn{3}{r}{\small Continúa en la página siguiente} \\
\endfoot

\hline
\endlastfoot

\texttt{3102.05} & Riego & Es el código principal del proyecto, ya que TierraViva se orienta a la automatización y supervisión del riego en cultivos desatendidos. La actuación física sobre la bomba se decide a partir de la humedad del suelo y de condiciones de seguridad definidas en el \textit{firmware}. \\
\hline
\texttt{3311.02} & Ingeniería de control & El prototipo incorpora una lógica de control local en cada estación: lectura de sensores, evaluación de umbrales, periodo de \textit{cooldown}, modo seguro y activación temporizada del relé. La estación no se limita a medir, sino que toma decisiones de actuación. \\
\hline
\texttt{3304.12} & Dispositivos de control & Las estaciones utilizan Arduino UNO como controlador local para coordinar sensores, comunicaciones, relé y bomba. El microcontrolador ejecuta la lógica de supervisión inmediata y garantiza que la actuación sobre el riego sea limitada y segura. \\
\hline
\texttt{3304.13} & Dispositivos de transmisión de datos & El sistema transmite telemetría desde las estaciones hacia ThingSpeak mediante el módulo A7670E-LASA y peticiones HTTP. La Raspberry Pi consulta posteriormente esos datos para construir el histórico y alimentar el \textit{dashboard}. \\
\hline
\texttt{3325.05} & Radiocomunicaciones & La conectividad de campo se apoya en comunicación móvil LTE. El módulo A7670E-LASA proporciona acceso remoto a Internet desde estaciones desplegadas en campo y permite publicar telemetría hacia ThingSpeak mediante red móvil. \\
\hline

\end{longtable}

\normalsize
\endgroup

\section{Códigos secundarios}

Los códigos secundarios describen áreas relevantes del desarrollo, aunque no constituyen por sí solas el objetivo final del proyecto. Se incluyen porque TierraViva integra \textit{firmware}, \textit{backend}, base de datos, API y visualización web.

\begingroup
\small
\renewcommand{\arraystretch}{1.18}
\rowcolors{2}{TableRowSoft}{white}

\begin{longtable}{|
>{\raggedright\arraybackslash}p{0.18\textwidth}|
>{\raggedright\arraybackslash}p{0.30\textwidth}|
>{\raggedright\arraybackslash}p{0.42\textwidth}|}

\caption{Códigos UNESCO secundarios vinculados al desarrollo software, sensorización y supervisión.}
\label{cuadro:codigos_unesco_secundarios}\\

\hline
\rowcolor{URJCRed}
\TVHead{Código UNESCO} & \TVHead{Denominación} & \TVHead{Justificación en TierraViva} \\
\hline
\endfirsthead

\hline
\rowcolor{URJCRed}
\TVHead{Código UNESCO} & \TVHead{Denominación} & \TVHead{Justificación en TierraViva} \\
\hline
\endhead

\hline
\multicolumn{3}{r}{\small Continúa en la página siguiente} \\
\endfoot

\hline
\endlastfoot

\texttt{1203.17} & Informática & El proyecto incluye desarrollo software en varias capas: \textit{firmware} de estación, \textit{scripts} de supervisión en Raspberry Pi, API FastAPI, almacenamiento SQLite y \textit{dashboard} web. \\
\hline
\texttt{1203.18} & Sistemas de información, diseño y componentes & La Raspberry Pi funciona como sistema de información local: consulta ThingSpeak, valida lecturas, almacena datos, expone \textit{endpoints} y proporciona información estructurada al \textit{dashboard}. \\
\hline
\texttt{1203.23} & Lenguajes de programación & TierraViva requiere programación en distintos entornos: C/C++ para Arduino, Python para \textit{backend} y supervisión, JavaScript para el \textit{dashboard} y \textit{scripts} auxiliares de despliegue. \\
\hline
\texttt{3304.17} & Sistemas en tiempo real & La estación ejecuta un ciclo periódico de lectura, decisión, actuación y publicación. La relación con este código se justifica por el control temporizado de la bomba, la gestión de retardos, el periodo de \textit{cooldown} y los tiempos máximos de actuación definidos en el \textit{firmware}. \\
\hline
\texttt{3307.03} & Diseño de circuitos & El proyecto incluye montaje electrónico, divisor de tensión, alimentación regulada, relé y cableado. El diseño se apoya en módulos comerciales integrados en una estación funcional, dejando el diseño de una placa electrónica propia como posible mejora futura. \\
\hline
\texttt{3311.07} & Instrumentos electrónicos & El prototipo utiliza sensores electrónicos para medir humedad del suelo, temperatura, humedad relativa, presión y luminosidad. Estas medidas constituyen la base de la monitorización y de la decisión local de riego. \\
\hline

\end{longtable}

\normalsize
\endgroup
